\documentclass[10pt,a4paper]{amsart}
\usepackage[margin=19mm]{geometry}
\usepackage{amsmath,amssymb,mathtools}
\usepackage[scr=rsfs,cal=euler]{mathalfa}
\usepackage{xfrac}
\usepackage[unicode,hidelinks,bookmarks=false]{hyperref}
\newcommand{\ie}{i.e.\ }
\newcommand{\cf}{cf.\ }
\newcommand{\resp}{respectively\ }
\newcommand{\Subsection}{\subsection}
\providecommand{\nobreakdash}{\nobreak\hbox{-}}
\setlength{\emergencystretch}{2em}
\multlinegap0pt
\theoremstyle{plain}
\newtheorem{theorem}{Theorem}
\newtheorem{lemma}{Lemma}
\newtheorem{proposition}{Proposition}
\newtheorem{corollary}[theorem]{Corollary}
\theoremstyle{remark}
\newtheorem{remark}{Remark}
\title[Chern flat manifolds that are torsion-critical]{Chern flat manifolds that are torsion-critical: every non-Kahler torsion-critical compact Chern flat manifold has semi-simple universal cover}
\author{Dongmei Zhang, Fangyang Zheng}
\date{Source: arXiv:2411.01132v1 (CC BY 4.0)}
\begin{document}
\maketitle
\noindent\emph{Source and licence.} Chern flat manifolds that are torsion-critical. Version arXiv:2411.01132v1 (CC BY 4.0), distributed under the Creative Commons Attribution 4.0 International licence, \url{http://creativecommons.org/licenses/by/4.0/}, as recorded in the arXiv licence metadata for this exact version and shown on the abstract page \url{https://arxiv.org/abs/2411.01132v1}. Full licence notice: licenses/arxiv-2411.01132-CC-BY-4.0.txt.\par\smallskip

\noindent\textit{Editorial scope.} Counted unit: the article's Proposition with source label \texttt{prop1} (source0.tex lines 168--170), the equivalence between the non-Kahler torsion-critical condition on a compact Chern flat manifold and the semi-simplicity of the universal cover, together with the converse existence of a compatible left-invariant metric on any semi-simple complex Lie group whose compact quotients are torsion-critical; delivered with the article's complete original proof of it (source0.tex lines 236--287, one \texttt{proof} environment of 52 lines whose heading names the counted proposition; the article places this proof after the auxiliary lemmas it uses). No mathematics is written, repaired or re-derived: statement and proof are the article's own text line for line, in the article's own notation (the Chern flat structure, the canonical metric, the compact real form, the torsion tensor and its components, the tensors A and B), and no retained line is substituted, re-flowed or omitted. Required context retained uncounted, in source order: the labelled display defining the torsion-critical condition (lines 138--140); the labelled display defining the tensors A and B (lines 147--149); and the statement of the article's first auxiliary lemma (lines 201--207), which the counted proof invokes in its final step. The proofs of those context results are not delivered. Nothing else of the article is retained: its other statements, its question and its remaining arguments are omitted. The counted proof cites ten of the article's references; the wrapper delivers exactly those ten entries, transcribed character for character from the article's own bibliography and carrying the article's own entry numbers as their printed labels, and no other bibliography entry is delivered and no bibliography database is shipped. Every cross-reference inside the retained spans resolves inside the delivered document. Editorial formatting only: an AMS \texttt{amsart} preamble that restates the environments the retained lines use, namely the article's declarations of Theorem, Lemma, Proposition, Corollary and Remark, and that adds no mathematical macro, because every control sequence used by the retained lines is standard. The article's own source declares the class \texttt{amsart} as well; no class or style file is shipped with this package. Numbering differs from the article: the retained auxiliary lemma prints as Lemma 1 and the counted proposition as Proposition 1, because the wrapper keeps the article's per-type counters without its sectional resets and retains only one environment of each type. No picture environment and no graphical inclusion occurs in any retained line, and the package delivers no image asset. One bundled source file, \texttt{sources/167845/source0.tex}, accompanies the delivered item as the exact source of every retained span. Every retained span is recorded in \texttt{delivery-evidence.json} with method \texttt{exact\_tex}.
\begin{equation} \label{eq:tor-crit}
2\sigma_1 - \sigma_2 +2(\phi + \bar{\phi}) - 2 (\xi + \bar{\xi}) = (|T|^2-\frac{n-1}{n}b)\, \omega, \ \ \ \ \ \ b=V^{-1}\int_M|T|^2dv.
\end{equation}

\begin{equation} \label{eq:AB}
A_{i\bar{j}} = \sum_{r,s}T^r_{is}\overline{T^r_{js}}, \ \ \ B_{i\bar{j}} = \sum_{r,s}T^j_{rs}\overline{T^i_{rs}}, \ \ \ \phi_{i}^j = \sum_r T^j_{ir} \overline{\eta}_r, \ \ \ \ \xi_{i}^j = \sum_r T^j_{ir,\,\bar{r}} ,
\end{equation}

\begin{proposition}  \label{prop1}
Let $(M^n,g)$ be a compact Chern flat manifold. Denote by $G$ its universal cover which is a complex Lie group. If $g$ is non-K\"ahler torsion-critical, then $G$ must be semi-simple. Conversely, given any semi-simple complex Lie group $G$, it admits a compatible left-invariant metric $g_0$ satisfying (\ref{eq:tor-crit}), so any compact quotient of $(G,g_0)$ is non-K\"ahler torsion-critical.  
\end{proposition}

\begin{lemma} \label{lemma1}
Let $(M^n,g)$ be a compact Chern flat manifold. Then it is torsion-critical if and only if 
\begin{equation} \label{eq:6}
2A - B = \frac{b}{n}I,
\end{equation} 
for some constant $b$. 
\end{lemma} 

\begin{proof}[{\bf Proof of Proposition \ref{prop1}.}]
We just need to show that any semi-simple complex Lie group $G$ admits a left-invariant metric compatible with its complex structure that is torsion-critical. Let $G$ be a  semi-simple complex Lie group. We want to show that its {\em canonical metric} $g_0$ on $G$ is torsion-critical. First let us  recall the definition of $g_0$ from \cite{PodestaZ}. The construction is well-known and we refer the readers to \cite{Gordon, Helgason, Lafuente} for more references. Here we will follow the notations and calculation in \cite{PodestaZ}.

Let ${\mathfrak g}$ be the Lie algebra of $G$, and we  denote by $\mathfrak g_{\mathbb R}$ the Lie algebra $\mathfrak g$ considered as a real algebra. A Cartan decomposition of  $\mathfrak g_{\mathbb R}$ is given by ${\mathfrak g_{\mathbb R}} = {\mathfrak u} + J{\mathfrak u}$, where ${\mathfrak u}$ is the Lie algebra of a compact semi-simple Lie group $U$. Let us denote by $K$ the {\em Cartan-Killing form} of ${\mathfrak g_{\mathbb R}}$, namely, 
$$ K(x,y) = \mbox{tr}(\mbox{ad}_x\mbox{ad}_y), \ \ \ \forall \ x,y \in {\mathfrak g_{\mathbb R}}. $$
The canonical metric $g_0$ on ${\mathfrak g_{\mathbb R}}$ is defined by
$$g_0|_{{\mathfrak u} \times {\mathfrak u}} = -K,\ \ \ g_0({\mathfrak u},J{\mathfrak u})=0,\ \ \ g_0|_{J{\mathfrak u} \times J{\mathfrak u}} = K.$$
Note that this metric depends on the choice of a compact real form ${\mathfrak u}$, but as two such real forms are conjugate by an inner automorphism of $\mathfrak g_{\mathbb R}$, the corresponding metrics are holomorphically isometric. 

Since $J$ is an integrable complex structure on ${\mathfrak g_{\mathbb R}}$, we have by definition that
$$ [x,y]-[Jx,Jy] + J[Jx,y] + J[x,Jy] = 0, \ \ \ \ \forall \ x,y \in {\mathfrak g_{\mathbb R}}. $$
Also,  ${\mathfrak g}$ being a complex Lie algebra means that 
$$ [x,y] + [Jx,Jy] = 0,  \ \ \ \ \forall \ x,y \in {\mathfrak g_{\mathbb R}}. $$
Combine the two we get
\begin{equation*}
[Jx,y]=J[x,y], \ \ \ \ \forall \ x,y \in {\mathfrak g}_{\mathbb R}. 
\end{equation*}
Now let $\{ u_1, \ldots , u_n, Ju_1, \ldots , Ju_n\}$ be an orthonormal basis of ${\mathfrak g}_{\mathbb R} = {\mathfrak u}+ J{\mathfrak u}$, and for convenience write $Ju_i=u_{i^{\ast}}=u_{n+i}$, $1\leq i\leq n$. 
Denote by $S^j_{ik}$ the structure constants of ${\mathfrak u}$, we have
$$ [u_i,u_k] =\sum_{j=1}^n S^j_{ik} u_j = - [u_{i^{\ast}}, u_{k^{\ast}}], \ \ \ [u_i, u_{k^{\ast}}] = \sum_{j=1}^n S^j_{ik}u_{j^{\ast}}. $$
For any $1\leq i,j,k \leq n$, we have
$$ \mbox{ad}_{u_i}\mbox{ad}_{u_k}(u_j) = [u_i, [u_k, u_j]] = \sum_{r,\ell =1}^n  S^{\ell}_{ir} S^r_{kj} u_{\ell}, \ \ \ \ \ \mbox{ad}_{u_i} \mbox{ad}_{u_k}(u_{j^{\ast}}) = \sum_{r,\ell =1}^n  S^{\ell}_{ir} S^r_{kj} u_{\ell^{\ast}}. $$
Therefore we have 
\begin{equation} \label{eq:SS}
\delta_{ik}=g_0(u_i,u_k) = - \mbox{tr}(\mbox{ad}_{u_i}\mbox{ad}_{u_k}) = -2 \sum_{j,r=1}^n  S^{j}_{ir} S^r_{kj}, \
 \ \ \ \forall \ 1\leq i,k\leq n.
\end{equation}

Next let us write $e_i=\frac{1}{\sqrt{2}}(u_i - \sqrt{-1}Ju_i)$, $1\leq i\leq n$. Then $\{ e_1, \ldots , e_n\}$ becomes a unitary basis for $({\mathfrak g}, g_0)$. Under the basis $e$, we compute
\begin{equation*}
[e_i,e_k] \,= \,\frac{1}{2}[ u_i - \sqrt{-1}Ju_i, \,u_k - \sqrt{-1}Ju_k] \, = \, [u_i,u_k] - \sqrt{-1}J [u_i,u_k] \, = \, \sqrt{2}\sum_{j=1}^n S^j_{ik}e_j.
\end{equation*}
Thus the Chern torsion components under $e$ becomes 
\begin{equation} \label{eq:TS}
T^j_{ik} = - \sqrt{2}S^j_{ik}, \ \ \ \ \ 1 \leq i,j,k\leq n. 
\end{equation}
It has been proved in \cite[Theorem 1.2]{PodestaZ} that $g_0$ is Bismut torsion parallel, namely, $\nabla^bT^b=0$ where $\nabla^b$ is the Bismut connection and $T^b$ is its torsion (see \cite{Bismut}. Since $T^b$ can be expressed in the Chern torsion $T^c=T$, this condition is equivalent to $\nabla^bT=0$. Bismut torsion parallel manifolds, and in particular those Hermitian manifolds where the Bismut curvature obeys all K\"ahler symmetries, were actively studied in recent years. See \cite{YZZ, ZZ, Zheng} and the references therein for more discussions). Hence $\nabla^bC=0$ where the symmetric $2$-tensor $C$ is defined by $C_{ij}=\sum_{r,s=1}^n T^r_{is}T^s_{jr}$ under any unitary frame. Under our particular choice of frame $e$ above, we have by (\ref{eq:SS}) and (\ref{eq:TS}) that
$$ C_{ij} = \sum_{r,s=1}^n T^r_{is}T^s_{jr} = 2 \sum_{r,s=1}^n S^r_{is}S^s_{jr} = -\delta_{ij}, \ \ \ \ \forall \ 1\leq i,j\leq n. $$
On the other hand, since $g_0$ is Chern flat, the Bismut connection $\nabla^b$ has coefficients (see for instance \cite{YZ} or \cite{VYZ})
$$ \nabla^b_{e_k}e_i = \sum_{j=1}^n T^j_{ik}e_j. $$
So the fact that $\nabla^bC=0$ yields
$$ \sum_{r=1}^n \big( T^r_{ik}C_{rj} + T^r_{jk}C_{ir}\big) = 0, \ \ \ \ \forall  \ 1\leq i,j,k\leq n. $$
That is, 
\begin{equation}
 T^j_{ik}  + T^i_{jk} =0 , \ \ \ \ \ \ \ \forall \ 1\leq i,j,k\leq n.
\end{equation}
Since $T$ is skew-symmetric with respect to its lower two indices, we also have $T^j_{ki}+T^i_{kj}=0$. By the definition (\ref{eq:AB}) for tensors $A$ and $B$, we have
$$ A_{i\bar{j}} =\sum_{r,s=1}^n T^r_{is} \overline{T^r_{js}} = - \sum_{r,s=1}^n T^s_{ir} \overline{T^r_{js}} = -2 \sum_{r,s=1}^n S^s_{ir} S^r_{js} =\delta_{ij}. $$
Similarly,
$$ B_{i\bar{j}} =\sum_{r,s=1}^n T^j_{rs} \overline{T^i_{rs}} = - \sum_{r,s=1}^n T^s_{jr} \overline{T^r_{is}} = -2 \sum_{r,s=1}^n S^s_{jr} S^r_{is} =\delta_{ij}. $$
So $A=B=I$ and by Lemma \ref{lemma1} we know that the canonical metric  $g_0$ on the semi-simple complex Lie group $G$ is torsion-critical. This completes the proof of Proposition \ref{prop1}. 
\end{proof}

\begin{thebibliography}{99}
\bibitem[2]{Bismut} J.-M. Bismut, \emph{A local index theorem for non-K\"ahler manifolds,} Math. Ann. {\bf 284} (1989), no.\,4, 681-699.
\bibitem[6]{Gordon} C. S. Gordon, \emph{Naturally reductive homogeneous Riemannian manifolds,} Can. J. Math. \textbf{37} (1985), 467-487.
\bibitem[7]{Helgason} S. Helgason, \emph{Differential Geometry, Lie groups and symmetric spaces,} Academic Press, Inc. (1978).
\bibitem[9]{Lafuente} R. A. Lafuente, M. Pujia, L. Vezzoni, \emph{Hermitian curvature flow on unimodular lie groups and static invariant metrics,} Trans. Amer. Math. Soc., \textbf{373} (2020), 3967-3993.
\bibitem[10]{PodestaZ} F. Podest\`a and F. Zheng, \emph{A note on compact homogeneous manifolds with Bismut parallel torsion,} arXiv: 2310.14002
\bibitem[11]{VYZ} L. Vezzoni, B. Yang, and F. Zheng, \emph{ Lie groups with flat Gauduchon connections,} Math. Zeit. \textbf{293} (2019), Issue 1-2, 597-608.
\bibitem[13]{YZ} B. Yang and F. Zheng, \emph{On curvature tensors of Hermitian manifolds,} Comm. Anal. Geom. {\bf 26} (2018), no.\,5, 1193-1220.
\bibitem[14]{YZZ} S.-T. Yau, Q. Zhao, and F. Zheng, \emph{On Strominger K\"ahler-like manifolds with degenerate torsion,} Trans. Amer. Math. Soc. \textbf{376} (2023), no.\,5, 3063-3085.
\bibitem[16]{ZZ} Q. Zhao and F. Zheng, \emph{Strominger connection and pluriclosed metrics,} J. Reine Angew. Math. (Crelles), \textbf{796} (2023), 245-267.
\bibitem[17]{Zheng} F. Zheng, \emph{Some recent progress in non-K\"ahler geometry,} Sci. China Math., {\bf 62} (2019), no.\,11, 2423-2434.
\end{thebibliography}
\end{document}
